\documentclass[11pt,a4paper]{article}
\usepackage{fontspec}
\usepackage{kotex}
\usepackage{geometry}
\usepackage{graphicx}
\usepackage{xcolor}
\usepackage{tcolorbox}
\usepackage{booktabs}
\usepackage{array}
\usepackage{longtable}
\usepackage{colortbl}
\usepackage{fancyhdr}
\usepackage{titlesec}
\usepackage{hyperref}
\usepackage{emoji}
\usepackage{amsmath}
\usepackage{amssymb}
\usepackage{enumitem}

\geometry{left=25mm, right=25mm, top=30mm, bottom=30mm}

\setmainfont{Noto Sans CJK KR}
\setsansfont{Noto Sans CJK KR}
\setmonofont{Noto Sans CJK KR}
\setemojifont{Noto Color Emoji}

\definecolor{primary}{HTML}{1976D2}
\definecolor{darkbrown}{HTML}{3E2723}
\definecolor{mediumbrown}{HTML}{5D4037}
\definecolor{lightbrown}{HTML}{8D6E63}
\definecolor{cream}{HTML}{FFF3E0}
\definecolor{lightgray}{HTML}{F5F5F5}
\definecolor{tableheader}{HTML}{E8E8E8}

\titleformat{\section}
  {\Large\bfseries\color{primary}}
  {\thesection.}{0.5em}{}
\titleformat{\subsection}
  {\large\bfseries\color{darkbrown}}
  {\thesubsection}{0.5em}{}
\titleformat{\subsubsection}
  {\normalsize\bfseries\color{mediumbrown}}
  {\thesubsubsection}{0.5em}{}

\renewcommand{\contentsname}{목차}
\renewcommand{\figurename}{그림}
\renewcommand{\tablename}{표}

\tcbset{
  tipbox/.style={colback=cream,colframe=primary,boxrule=1.5pt,arc=3mm,
    left=5mm,right=5mm,top=3mm,bottom=3mm,fonttitle=\bfseries\color{primary}},
  formulabox/.style={colback=lightgray,colframe=lightgray,boxrule=0pt,arc=2mm,
    left=5mm,right=5mm,top=3mm,bottom=3mm},
  summarybox/.style={colback=white,colframe=primary,boxrule=2pt,arc=4mm,
    left=5mm,right=5mm,top=5mm,bottom=5mm}
}

\hypersetup{colorlinks=true,linkcolor=primary,urlcolor=primary,bookmarks=true,unicode=true}


\pagestyle{fancy}
\fancyhf{}
\fancyhead[R]{\small\color{lightbrown}MakeCode LED 이름 표시 시뮬레이터 이론해설서 v1.0.0.0}
\fancyfoot[C]{\thepage}
\renewcommand{\headrulewidth}{0.4pt}

\begin{document}

\begin{titlepage}
\centering
\vspace*{3cm}
{\Huge\bfseries\color{primary}\emoji{desktop-computer} MakeCode LED 이름 표시 시뮬레이터}\\[0.5cm]
{\Huge\bfseries\color{primary} 이론해설서}
\\[1.5cm]
{\Large 시뮬레이터 버전: v1.0.2.0}
\\[2cm]
\begin{tcolorbox}[summarybox]
\centering
{\large 이 이론해설서는 MakeCode LED 이름 표시 시뮬레이터이 다루는 핵심 개념과 원리를 설명합니다.}
\end{tcolorbox}
\vfill
{\large 사물인터넷과 빅데이터}\\[0.3cm]
{\large 청주교육대학교 컴퓨터교육과}
\vspace{1cm}
\end{titlepage}

\section*{1. 이 문서의 목적}

이 이론해설서는 \textbf{MakeCode LED 이름 표시 시뮬레이터}이 다루는 핵심 개념과 원리를 설명합니다. 시뮬레이터를 사용하기 전에 배경 지식을 쌓거나, 사용 후 깊이 있는 이해를 위해 참조합니다.


\begin{tcolorbox}[summarybox]
\textbf{핵심 주제}: 블록코딩과 마이크로비트

\end{tcolorbox}


\section*{2. 핵심 개념}

\begin{center}
\renewcommand{\arraystretch}{1.4}
\begin{tabular}{|p{3.5cm}|p{11.5cm}|}
\hline
\rowcolor{tableheader}
\textbf{개념} & \textbf{설명} \\
\hline
\textbf{블록코딩} & 텍스트 대신 시각적 블록을 조합하여 프로그래밍하는 방식. 문법 오류 없이 논리에 집중 가능 \\
\hline
\textbf{LED 매트릭스} & 5\texttimes{}5 = 25개 LED가 격자로 배열된 디스플레이. 각 LED를 개별 제어 가능 \\
\hline
\textbf{이벤트 기반 프로그래밍} & 특정 사건(버튼 클릭, 센서 값 변화 등)이 발생할 때 코드가 실행되는 방식 \\
\hline
\textbf{순차 실행} & 블록이 위에서 아래로 순서대로 실행되는 프로그래밍의 기본 원리 \\
\hline
\textbf{마이크로비트} & 영국 BBC가 개발한 교육용 마이크로컨트롤러. ARM Cortex-M4, BLE, 센서 내장 \\
\hline
\end{tabular}
\end{center}


\section*{3. 원리 상세}

\subsection*{블록코딩의 교육적 가치}

블록코딩은 MIT 미디어랩의 스크래치(Scratch)에서 시작된 접근으로, 프로그래밍의 본질인 \textbf{논리적 사고}에 집중할 수 있게 합니다.


\textbf{텍스트 코딩 vs 블록코딩:}

\begin{center}
\renewcommand{\arraystretch}{1.4}
\begin{tabular}{|p{2.5cm}|p{6.8cm}|p{5.7cm}|}
\hline
\rowcolor{tableheader}
\textbf{구분} & \textbf{텍스트 코딩} & \textbf{블록코딩} \\
\hline
문법 오류 & 잦음 (세미콜론, 괄호 등) & 없음 (블록이 맞춰짐) \\
\hline
진입 장벽 & 높음 & 낮음 \\
\hline
표현력 & 매우 높음 & 제한적 \\
\hline
적합 대상 & 전문 개발자 & 초보자, 교육 환경 \\
\hline
\end{tabular}
\end{center}

\subsection*{MakeCode와 마이크로비트}

Microsoft MakeCode는 블록↔JavaScript 양방향 변환을 지원하는 교육용 코딩 환경입니다. 마이크로비트의 LED, 버튼, 센서를 블록으로 쉽게 제어할 수 있습니다.


\textbf{이벤트 기반 프로그래밍:}

$\bullet$ \texttt{시작하면 실행}: 전원이 켜질 때 1번 실행

$\bullet$ \texttt{버튼 A 누르면}: 사용자 입력에 반응

$\bullet$ \texttt{흔들면}: 가속도계 센서가 감지


이 패턴은 스마트폰 앱, 웹사이트, 게임 등 모든 인터랙티브 소프트웨어의 기본입니다.



\section*{4. 실생활 응용 사례}

\textbf{사례 1}: 스마트워치: LED/화면에 시간·알림을 표시하는 원리


\textbf{사례 2}: 전광판: LED 매트릭스로 텍스트·이미지를 스크롤 표시


\textbf{사례 3}: 신호등: 이벤트 기반으로 빨강\textrightarrow{}노랑\textrightarrow{}초록 순서 제어


\textbf{사례 4}: 엘리베이터: 버튼 입력 \textrightarrow{} 층 이동 \textrightarrow{} LED 표시의 순차 실행



\section*{5. 더 알아보기}

$\bullet$ \textbf{MicroPython}: 마이크로비트에서 텍스트 코딩을 할 수 있는 Python 변형

$\bullet$ \textbf{GPIO 핀}: LED 매트릭스 외에 외부 장치를 제어하는 핀

$\bullet$ \textbf{인터럽트}: 이벤트 기반 프로그래밍의 하드웨어 레벨 원리

$\bullet$ \textbf{임베디드 시스템}: 특정 기능을 수행하도록 설계된 컴퓨터 시스템



\section*{6. 핵심 용어 사전}

\begin{center}
\renewcommand{\arraystretch}{1.4}
\begin{tabular}{|p{3.3cm}|p{1.8cm}|p{9.8cm}|}
\hline
\rowcolor{tableheader}
\textbf{용어} & \textbf{학술 용어} & \textbf{쉬운 설명} \\
\hline
블록코딩 & — & 텍스트 대신 시각적 블록을 조합하여 프로그래밍하는 방식 \\
\hline
LED 매트릭스 & LED 매트릭스 & 5\texttimes{}5 = 25개 LED가 격자로 배열된 디스플레이 \\
\hline
이벤트 기반 프로그래밍 & — & 특정 사건(버튼 클릭, 센서 값 변화 등)이 발생할 때 코드가 실행되는 \\
\hline
순차 실행 & — & 블록이 위에서 아래로 순서대로 실행되는 프로그래밍의 기본 원리 \\
\hline
마이크로비트 & — & 영국 BBC가 개발한 교육용 마이크로컨트롤러 \\
\hline
\end{tabular}
\end{center}


\section*{7. 참고자료}

\begin{center}
\renewcommand{\arraystretch}{1.4}
\begin{tabular}{|p{1.5cm}|p{7.6cm}|p{5.9cm}|}
\hline
\rowcolor{tableheader}
\textbf{\#} & \textbf{자료명} & \textbf{비고} \\
\hline
1 & MakeCode시뮬레이터LED이름표시\_퀵가이드 & 소프트웨어 사용 중 빠른 참조 \\
\hline
2 & MakeCode시뮬레이터LED이름표시\_사용설명서 & 전체 기능 상세 \\
\hline
3 & MakeCode시뮬레이터LED이름표시\_활용가이드 & 수업 적용 \\
\hline
4 & MakeCode시뮬레이터LED이름표시\_개발참조 & 기술 사양 \\
\hline
\end{tabular}
\end{center}


\vfill
\begin{center}
\small\color{lightbrown}
\textit{MakeCode LED 이름 표시 시뮬레이터 이론해설서}
\end{center}
\end{document}